\originalchapter{来源映射、范围与缺项}
\originalsection{18份课堂讲义的逐文件映射}
原作者均为 Jared Speck, 原课程 MIT 18.152, Fall 2011. 表中页数是实际 PDF 页序, 各含末尾一页 OCW 来源提示, 不是估计正文页数. 所有18份文件哈希均不同, 无整份重复; 合并的是重复概念和论证, 不删除课程主题. 文件名前缀均为 \texttt{mit18\_152f11\_lec\_}. 原件链接及完整 SHA256 见本目录 \texttt{source-manifest.json}, 原件只位于仓库 sources.
\begin{longtable}{p{.13\textwidth}p{.08\textwidth}p{.35\textwidth}p{.32\textwidth}}
\toprule 文件尾号&页数&原讲次及主题&中文正文位置\\\midrule\endhead
01&8&1: PDE与线性性&第1、13章\\
02&8&2: 热方程与分离变量&第2、3章\\
03&4&3: 热唯一性&第2章\\
04&4&4: 热最大原理&第2、4章\\
05&11&5: 热核及全空间&第4章\\
06&7&6: Laplace/Poisson、平均值&第5、6章\\
07&11&7: Poisson 基本解&第6章\\
08&4&8: Green 函数&第7章\\
09&7&9: Poisson 核、Harnack、Liouville&第5、7章\\
10&7&10: 波导论、一维传播&第8章\\
11&4&11: 球面平均&第8章\\
12&7&12: Kirchhoff、Minkowski&第8、9章\\
13\_14&7&13--14: 几何能量、应力张量&第9章及作业/Bonus\\
15&8&15: 二阶方程分类&第10章\\
16\_18&12&16--18: Fourier、反演、Plancherel&第10章\\
19\_20&10&19--20: Schrödinger&第11章\\
21\_23&11&21--23: Lagrange 场论&第12章及作业11\\
24&6&24: 输运与 Burgers&第13章\\\bottomrule
\end{longtable}
\originalsection{课程考核完整性与答案身份}
官方公开11份作业、1份Bonus、期中卷、期中官方答案和期末卷, 共15个PDF、78页. 大量考试空白答题页不算题目. 11份作业有51个罗马编号条目, 其中39道完整公开大题、11条仅商业教材引用、1条阅读任务. 加Bonus1、期中5、期末7, 收入52道完整公开大题、114个叶条目. 无字母分问的大题按1叶计, 内部多个要求仍全部解答. 三个完全重复小问复用已写证明, 不删除其课程位置; 其余111叶各有独立任务.

作业1--11依序对应原题 I--IV; III--V; I、II、V、VI; VI; I--III; I--III; I--V; I、II; I--VI; I--IV; I--IV. 叶条目数依次为4、3、4、1、3、7、7、7、13、13、13. Bonus原卷两个f保留为 $f_1,f_2$, 共10叶. 期中 I--V共9叶, 期末 I--VII共20叶. 每一道题面注明原题号, 可通过目录进入相应作业或考试. Bonus(a,b,$f_1$)分别完全重复作业8-II(a,b,e), 在Bonus明确回指.

只有期中有官方公开答案, 本册基于官方答案重构并明确条件修订. 其余作业、Bonus、期末的答案为 AI 独立编写及交叉数学审读, 不署为MIT官方答案. 官方作业8的错误迹/守恒式由官方Bonus修正题面交叉核验; Bonus(h)点态等式只在积分后成立. 其余发现的原件误号、符号、初值角点和解类限制在对应题处明示, 不暗示官方已发布勘误. 独立审查报告和逐题覆盖记录位于 \texttt{review/}.
\originalsection{商业教材引用的真实缺项}
下表仅登记官方公开作业中的引用, 未获得完整题面, 不猜造问题、不声称已完成解答. 指定书为 Salsa, \textit{Partial Differential Equations in Action: From Modelling to Theory}, Springer, 2010.
\begin{longtable}{p{.2\textwidth}p{.2\textwidth}p{.12\textwidth}p{.35\textwidth}}
\toprule 作业原号&书中题号&书页&公开补充信息\\\midrule\endhead
2-I&2.1&97&无完整题面\\
2-II&2.3&97&仅给 $L=\pi,U=0$\\
3-III&2.13&99&无完整题面\\
3-IV&2.14&99&允许热核表示的提示\\
4-I&3.1&150&无完整题面\\
4-II&3.2&150&无完整题面\\
4-III&3.3&151&补 $u\in C^3(\Omega)\cap C^1(\overline\Omega)$\\
4-IV&3.4&151&无完整题面\\
4-V&3.8&152&无完整题面\\
5-IV&3.14&153&Green 负性与对称性提示, 非题面\\
5-V&3.21&155&Helmholtz 基本解提示, 非题面\\\bottomrule
\end{longtable}
作业1-V要求阅读书的附录A, 只登记阅读任务. 这些商业书引用不受本项目的OCW资源许可自动授权, 源码包不含该书或猜造的书题.
\originalsection{核心范围与编者补充}
本册覆盖全部18份公开课堂讲义的数学主线: 热初边值、能量与最大原理、Gaussian 基本解及增长受控的全空间问题; 调和平均值、Poisson基本解与Green核、Poisson公式、Harnack和Liouville; 一维和三维显式波传播、二维下降法、锥能量与有限传播; 主部分类、Fourier反演与Plancherel; 自由Schrödinger及色散; Lorentz应力张量、变分原理、等距流与守恒; 线性输运和Burgers首次破裂. 核心结论给证明, 公开课程题逐项中文作答.

编者补充包括: 分布初值和基本解的测试函数表述; 热级数的 $L^2$ 与角点解释; Fejér完备性证明; Poisson能量弱解的Riesz构造; Duhamel换积分条件; Fourier频率截断的能量类唯一性; 含势边界能量条件; Burgers弱激波、Rankine--Hugoniot与基本熵Riemann解. 这些帮助连接已有实分析、泛函和ODE基础, 不把它们标作原课额外公开讲义.

未建立的范围: 一般粗糙域椭圆边界正则性与全套Sobolev理论; 一般非线性PDE的局部/全局存在; 任意Burgers数据的熵解完整存在唯一性; 任意势的自伴算子定义域; Maxwell系统、流体力学系统及高维非线性色散理论; 有限差分、有限元与数值收敛. 原课程商业教材全书及11处题面缺项也不在已覆盖范围.
\originalsection{署名、许可与变更}
原课程材料: Jared Speck, MIT OpenCourseWare, 18.152 Introduction to Partial Differential Equations, Fall 2011. 中文整理为翻译、重排、增补证明、AI题解、条件修订和自绘图. 改编采用 CC BY-NC-SA 4.0, 保留署名、非商业和相同许可要求, 不表示原作者或MIT认可. 逐个PDF核对实际页数、首页作者、末页OCW提示、资源页许可及第三方声明: 共33份, 未检出单独第三方版权声明. 商业教材是明确的非公开来源边界. 权利结论依据归档原件及当时的官方条款, 不概括为所有MIT材料均无例外.

官方入口: \href{https://ocw.mit.edu/courses/18-152-introduction-to-partial-differential-equations-fall-2011/pages/lecture-notes/}{课堂讲义}, \href{https://ocw.mit.edu/courses/18-152-introduction-to-partial-differential-equations-fall-2011/pages/assignments/}{作业与Bonus}, \href{https://ocw.mit.edu/courses/18-152-introduction-to-partial-differential-equations-fall-2011/pages/exams/}{考试和期中答案}, \href{https://ocw.mit.edu/pages/privacy-and-terms-of-use/}{OCW许可条款}.
